%======================================================================
\section{The ladder below $\Gamma$}\label{sec:ladder}
%======================================================================

Before $\Gamma$ we found two groups that each supply part of what the criterion needs. We record them here for comparison, and because each
leaves an open problem; nothing in this section is used in the proof of Theorem~A or in Part~II. The dyadic group $\mathcal G$ of
Section~\ref{sec:dyadic} is elementary amenable and transports copies of $S_3$ by the powers of one element, with far commutators of area at most
quadratic. Corollary~\ref{cor:cyclic} turns this into the exponent $1/3$, the most that transport along a ray of $H_3$ can give
(Section~\ref{sec:criterion}). The subgroup $\mathcal L$ of Thompson's group $V$ in Section~\ref{sec:thompson} transports along the exponentially
growing orbits of Thompson's group $F$: words of length $O(\log N)$ reach $N$ sites, and the quadratic Dehn function of $F$ fills all pairwise
commutators cheaply. That comes within a factor $\log^2r$ in the exponent of the ceiling of Theorem~\ref{thm:criterion}, but inside a group that
contains $F$, whose soficity is unknown. The proof that $\mathcal L$ is hyperlinear exactly when $F$ is uses tools that nothing else in the paper
needs, and it is in Appendix~\ref{app:hyperlinear}.

\subsection{An amenable group at exponent \texorpdfstring{$1/3$}{1/3}}\label{sec:dyadic}

The group $\mathcal G$ is a lamplighter over the dyadic rationals $\mathbb Z[\frac12]$: the lamps are copies of $S_3$, and the lighter is
generated by the shift of the integers and by doubling. In this section a group acting on a set $X$ acts on $\bigoplus_{x\in X}S_3$ by carrying
the copy of $S_3$ at $x$ to the copy at $g(x)$ through the identity of $S_3$.

Let $X=\mathbb Z[\frac12]$, and let $t$ and $u$ be the permutations of $X$ given by $t(x)=x+1$ for $x\in\mathbb Z$, $t(x)=x$ for
$x\notin\mathbb Z$, and $u(x)=2x$. Let $\mathcal B=\langle t,u\rangle$ and $\mathcal G=\bigl(\bigoplus_{x\in X}S_3\bigr)\rtimes\mathcal B$, and let $a,b$ be
generating involutions of the copy of $S_3$ at $0$. Put $s=utu^{-1}$ and $v=tst^{-1}$. Then $s$ adds $2$ to the even integers and $v$ adds $2$
to the odd integers, each fixing every other point, so $t^2=sv$ and $[s,v]=1$. Moreover $u$ and $v$ fix $0$, and $s$ fixes $1=t(0)$. Since
$\mathcal B$ acts transitively on $X$, the group $\mathcal G$ is generated by $a$, $b$, $t$ and $u$. By these identities the following seventeen words
in $a,b,t,u,s,v$ represent $1$ in $\mathcal G$; they form the set $R_{\mathcal G}$:
\[
a^2,\ b^2,\ (ab)^3,\quad s^{-1}utu^{-1},\ v^{-1}tst^{-1},\ t^2v^{-1}s^{-1},\ [s,v],\quad [u,x],\ [v,x],\ [s,txt^{-1}],\ [x,tyt^{-1}],
\]
where $x,y\in\{a,b\}$. Theorem~\ref{thm:dyadic} also records that $\mathcal G$ is not solvable, so the exponent $1/3$ is attained outside the
solvable class.

\begin{theorem}[The dyadic group]\label{thm:dyadic}
The group $\mathcal G$ is elementary amenable and not solvable, and every commutator $[x,t^kyt^{-k}]$ with $x,y\in\{a,b\}$ and
$k\ge1$ has area at most $5k^2$ with respect to $R_{\mathcal G}$. Consequently the chunk $E_{\mathcal G}$ that Theorem~\ref{thm:criterion}
associates with $R_{\mathcal G}$ has finite sofic profile with
\[
\log\sigma_{E_{\mathcal G}}(r)\ \ge\ \frac1{16}\Big\lfloor\Big(\frac r{75\cdot10^8}\Big)^{1/3}\Big\rfloor\qquad(r>1),
\]
every unitary assignment $\phi$ as in Theorem~\ref{thm:criterion} has $\log D\ge\frac{\Delta^2}{12}\bigl((5K_\Delta e)^{-2/5}-1\bigr)$, and
$\mathcal G$ is not LEF.
\end{theorem}

\begin{proof}
\emph{Amenability.} For $j\in\mathbb Z$ let $t_j=u^jtu^{-j}$; it adds $2^j$ to the points of $2^j\mathbb Z$ and fixes the others. The group $\mathcal K$
generated by all $t_j$ is normalized by $u$ and contains $t$, so it is normal in $\mathcal B$ and $\mathcal B/\mathcal K$ is cyclic. For $n\ge1$ the elements $t_j$ with
$|j|\le n$ move only points of $2^{-n}\mathbb Z$, and they commute with the translation $x\mapsto x+2^n$ of $2^{-n}\mathbb Z$, which has $4^n$
orbits. A permutation of $2^{-n}\mathbb Z$ that commutes with this translation maps each orbit onto an orbit by a translation, so these $t_j$ generate a subgroup of
$\mathbb Z^{4^n}\rtimes\Sym_{4^n}$, which is virtually abelian. Hence $\mathcal K$ is a directed union of virtually abelian groups, $\mathcal B$ is elementary
amenable, and so is $\mathcal G$, because $\bigoplus_XS_3$ is locally finite.

\emph{Not solvable.} For $n\ge1$ the elements $t_0,\dots,t_{n-1}$ preserve $\mathbb Z$ and commute with $x\mapsto x+2^n$, so they act on
$\mathbb Z/2^n$, where they permute the classes modulo $2^i$ for every $i\le n$. So they generate a subgroup $\Gamma_n$ of the iterated wreath
product $\mathsf W_n$ of $n$ copies of $\mathbb Z/2$, which has order $2^{2^n-1}$. On the even classes, identified with $\mathbb Z/2^{n-1}$ by
halving, $t_1,\dots,t_{n-1}$ act as $t_0,\dots,t_{n-2}$ act on $\mathbb Z/2^{n-1}$, and they fix the odd classes. Conjugation by $t_0$ carries
this action to the odd classes, and $t_0$ exchanges the two halves. So $|\Gamma_n|\ge2|\Gamma_{n-1}|^2$, and $\Gamma_n=\mathsf W_n$ by induction from $\Gamma_1=\mathsf W_1=\mathbb Z/2$. The derived
subgroup of $\mathsf W_n$ lies in $\mathsf W_{n-1}\times\mathsf W_{n-1}$ and projects onto each factor, since the commutator of $(g,1)$ with the
exchange is $(g^{-1},g)$. By induction the $(n-1)$st derived subgroup of $\mathsf W_n$ is nontrivial. Since $\Gamma_n$ is a quotient of the subgroup
$\langle t_0,\dots,t_{n-1}\rangle$ of $\mathcal B$, the derived length of $\mathcal B$ is at least $n$ for every $n$, so neither $\mathcal B$ nor
$\mathcal G$ is solvable.

\emph{Area.} Let $A_{\mathcal G}(k)$ be the largest area of $[x,t^kyt^{-k}]$ over $x,y\in\{a,b\}$, so that $A_{\mathcal G}(1)=1$. Five kinds of move
cost one cell each.
\begin{itemize}[nosep,leftmargin=*]
\item Replacing $t^2$ by $sv$.
\item Exchanging adjacent letters $s^{\pm1}$ and $v^{\pm1}$.
\item Replacing $s$ by $utu^{-1}$.
\item Moving a letter $x^{\pm1}$, $x\in\{a,b\}$, across $u^{\pm1}$ or $v^{\pm1}$.
\item Moving $txt^{-1}$ or its inverse across $s^{\pm1}$.
\end{itemize}
The commutator contains four powers $t^{\pm k}$.

For $k=2m$, write each $t^{\pm2m}$ as $(s^mv^m)^{\pm1}$, with $m$ replacements and $m(m-1)/2$ exchanges each, $2m^2+2m$ cells in all. Moving
$y^{\pm1}$ across $v^{\pm m}$ ($2m$ cells) leaves $[x,s^mys^{-m}]$. Replacing its $4m$ letters $s^{\pm1}$ gives
$[x,ut^mu^{-1}yut^{-m}u^{-1}]$, and moving $y^{\pm1}$ and $x^{\pm1}$ across $u$ (four cells) gives $u[x,t^myt^{-m}]u^{-1}$. Hence
$A_{\mathcal G}(2m)\le A_{\mathcal G}(m)+2m^2+8m+4$.

For $k=2m+1$, put $z=tyt^{-1}$, so that the commutator is freely $[x,t^{2m}zt^{-2m}]$. Write each $t^{\pm2m}$ as $(v^ms^m)^{\pm1}$, now with
$m(m+1)/2$ exchanges each, $2m^2+6m$ cells in all. Moving $z^{\pm1}$ across $s^{\pm m}$ ($2m$ cells) leaves $[x,v^mzv^{-m}]$, moving
$x^{\pm1}$ across $v^{\pm m}$ ($2m$ cells) leaves $v^m[x,z]v^{-m}$, and $[x,z]$ is a relation. Hence $A_{\mathcal G}(2m+1)\le2m^2+10m+1$.

By induction $A_{\mathcal G}(k)\le5k^2$, since $A_{\mathcal G}(2)\le15$, $A_{\mathcal G}(m)+2m^2+8m+4\le7m^2+8m+4\le20m^2$ for $m\ge1$, and
$2m^2+10m+1\le5(2m+1)^2$.

\emph{Profile.} The words of $R_{\mathcal G}$ have length at most $8$, so Corollary~\ref{cor:cyclic} with $C=5$, $p=2$ and $l=8$ gives the
bounds. Since $\mathcal G$ is amenable, it is sofic and every chunk has finite profile. The profile of $E_{\mathcal G}$ is unbounded, so
$E_{\mathcal G}$ embeds in no finite group, and $\mathcal G$ is not LEF (Section~\ref{sec:criterion}).
\end{proof}

The contrast with $H_3$ is instructive. There the doubling endomorphism of~\cite[Section~3]{DouglasMghirbiA}, which differs by a conjugation
from the doubling map of Burillo, Cleary, Martino and R\"over~\cite[Section~7]{BCMR2016}, multiplies a reweighted area by $107$ at each step;
hence the exponent $\log107$. In $\mathcal G$ the dilation $u$ relates distance $2m$ to distance $m$ exactly, and the two commuting shifts cost
$O(m^2)$ cells. The orbit of $0$ under $\mathcal B$ grows exponentially, as the orbits of Thompson's group $F$ do. Whether $\mathcal G$ also
supports the transport of many copies that we use for $F$ in Section~\ref{sec:thompson} is open (Section~\ref{sec:open}).

\subsection{Thompson's group \texorpdfstring{$V$}{V}}\label{sec:thompson}

R\"over proved that every Houghton group embeds in Thompson's group $V$~\cite[Proposition~2.6]{Rover1999}, and Genevois extended the
embedding to generalised Houghton groups~\cite[Corollary~5.11]{Genevois2019}. Since $H_3$ is amenable, Thompson's group $V$ therefore has a
chunk whose profile is finite at every $r$ and superpolynomial by~\cite[Theorem~1.1]{DouglasMghirbiA}. A different subgroup of $V$ gives a much
larger bound.

Let $\mathbb D$ be the set of dyadic rationals in $(0,1)$, and let Thompson's group $F$ act on $[0,1]$ by its piecewise-linear
homeomorphisms. Let $\mathcal L=\bigl(\bigoplus_{x\in\mathbb D}S_3\bigr)\rtimes F$. To embed $\mathcal L$ in $V$ we double every $1$ of a binary
word, so that the pattern $10$ never occurs inside a coded word. The first $10$ then marks a gap, the gaps are indexed by the dyadic rationals,
and $F$ acts on the gaps as it acts on $\mathbb D$.

\begin{lemma}[Gap coding]\label{lem:gap-coding}
$\mathcal L$ embeds in $V$.
\end{lemma}

\begin{proof}
Let $\mathsf c$ be the substitution $0\mapsto0$, $1\mapsto11$ on finite binary words. An infinite binary word either is a concatenation of
blocks $0$ and $11$, or it has a first occurrence of $10$ at a block boundary. So Cantor space is the disjoint union of the set $\mathcal C$ of
block sequences and the cylinders $Z_w=\mathsf c(w)10\{0,1\}^{\mathbb N}$, one for each finite word $w$. Index $Z_w$ by the dyadic
$x_w=0.w1\in\mathbb D$ in binary; this is a bijection. Let $f\in F$ have a tree-pair diagram with leaves $w_1,\dots,w_n$ and $w'_1,\dots,w'_n$,
so that $f(0.w_kz)=0.w'_kz$. Define $\hat f$ by the prefix replacements $\mathsf c(w_k)\mapsto\mathsf c(w'_k)$ and
$\mathsf c(v)10\mapsto\mathsf c(v')10$, where the interior vertices $v$ of the first tree are matched in left-to-right order, the $k$-th with the $k$-th, with the interior vertices $v'$ of
the second. Both families of prefixes are complete prefix codes, because splitting a leaf $w$ replaces $\mathsf c(w)$ by $\mathsf c(w0)$,
$\mathsf c(w)10$ and $\mathsf c(w1)$. Expanding a matched pair of leaves adds a matched pair of gaps and changes nothing else, so $\hat f$ depends
only on $f$. It maps $Z_w$ onto $Z_{w'}$ with $x_{w'}=f(x_w)$: below a leaf this is the formula for $f$, and interior vertices correspond to the
breakpoints between leaves. On $\mathcal C$ it acts as $f$ acts on binary sequences, so $f\mapsto\hat f$ is an injective homomorphism $F\to V$.
Inside each $Z_w$ let $S_3$ permute the three cylinders $\mathsf c(w)10y\{0,1\}^{\mathbb N}$, $y\in\{0,10,11\}$. Copies with different $w$ have
disjoint supports, and $\hat f$ carries the copy at $x$ to the copy at $f(x)$. An element of the group they generate that acts trivially acts
trivially on $\mathcal C$, so its $F$-part is trivial, and then each of its $S_3$ coordinates is trivial.
\end{proof}

Let $a,b$ generate the copy of $S_3$ at $\frac12$, let $t\in F$ be
\[
t(x)=\begin{cases}2x,&0\le x\le\frac14,\\ x+\frac14,&\frac14\le x\le\frac12,\\ \frac{x+1}2,&\frac12\le x\le1,\end{cases}
\]
so that $t(\frac12)=\frac34$, and for $m\ge1$ let $\mathbb D_m$ be the set of the $2^m-1$ points of $\mathbb D$ with denominator at most
$2^m$.

The group $\mathcal L$ is finitely presented. This is Cornulier's criterion for permutational wreath products~\cite[Theorem~1.1]{Cornulier2006},
applied to the action of $F$ on $\mathbb D$; he checks its hypotheses in~\cite[Example~3.4]{Cornulier2006}. Thompson's group $F$ is finitely
presented, and it acts transitively on increasing pairs of points of $\mathbb D$~\cite[Lemma~4.2]{CFP1996}, so it has finitely many orbits on
$\mathbb D\times\mathbb D$. The stabilizer $F_{1/2}$ of $\frac12$ is the product of the elements supported in $[0,\frac12]$ and those supported
in $[\frac12,1]$, two copies of $F$, so it is finitely generated. Nothing in Proposition~\ref{prop:thompson} uses finite presentation.

\begin{proposition}[Thompson transport]\label{prop:thompson}
The group $\mathcal L$ contains $F$ and contains no Houghton group $H_m$, so the bounds below do not come from the
Houghton subgroups of $V$. There are a finite set $R$ of relations of $\mathcal L$, containing $a^2,b^2,(ab)^3$, and a constant $C$ with the
following property: for every $m\ge1$ there are words $w_p$, $p\in\mathbb D_m$, with $w_p(\frac12)=p$ and
\[
\mathrm{Area}_R\bigl([w_pxw_p^{-1},w_qyw_q^{-1}]\bigr)\le Cm^2\qquad(p\ne q\in\mathbb D_m,\ x,y\in\{a,b\}).
\]
Consequently a chunk $E$ of $\mathcal L$, and hence of $V$, has $\log\sigma_E(r)\ge c\,r/\log^2r$ for all large $r$, and for $0<\Delta\le1$
every unitary assignment as in Theorem~\ref{thm:criterion} has $\log D\ge c_\Delta\,e^{-2}/\log^4(1/e)$ for all small $e$, where $c>0$ is a
constant and $c_\Delta>0$ depends only on $\Delta$.
\end{proposition}

\begin{proof}
\emph{Transport.} A standard dyadic interval is one of the form $[k2^{-j},(k+1)2^{-j}]$. Let two partitions of $[0,1]$ into standard dyadic
intervals have the same number $n$ of intervals. The map sending the $i$th interval of the first affinely onto the $i$th interval of the second
lies in $F$ and has a tree-pair diagram with $n-1$ carets. Every standard dyadic interval can be split into any given number of standard
dyadic intervals. For $p<q$ in $\mathbb D_m$, the binary digits of $p$ split $[0,p]$ and $[p,1]$ into at most $m$ standard dyadic
intervals each. The interval $[p,q]$, cut at its point $c$ with the shortest binary expansion, splits into at most $m$ standard dyadic intervals
on each side of $c$. Hence there are $w_p\in F$ with $w_p(\frac12)=p$ and fewer than $2m$ carets, and $g_{pq}\in F$ with $g_{pq}(\frac12)=p$, $g_{pq}(\frac34)=q$ and fewer than $4m$
carets. The word length in $F$ is comparable to the number of carets of the reduced tree-pair diagram~\cite{BCS2001}. Hence these elements
are represented by words of length $O(m)$ in the generating set of the finite presentation of $F$ that the set $R$ of relations uses. The Dehn
functions of two finite presentations of $F$ are equivalent, so Guba's quadratic bound holds for it. We write $w_p$ also for such a word. An
element fixing $\frac12$ maps $[0,\frac12]$ to itself, so the left and right subtrees of its reduced diagram are diagrams for its two
restrictions, which are elements of $F$ after affine rescaling, with fewer carets in total. Hence $F_{1/2}\cong F\times F$ is undistorted in
$F$: its elements of length $n$ in $F$ are words of length $O(n)$ in a fixed finite set of generators $h_j$ of $F_{1/2}$.

\emph{Relations.} Let $R$ consist of a finite presentation of $F$ on a generating set that includes $t$ and the $h_j$, the words
$a^2,b^2,(ab)^3$, the commutators $[x,h_j]$, and the commutators $[x,tyt^{-1}]$, for $x,y\in\{a,b\}$. They hold in $\mathcal L$: each $h_j$
fixes $\frac12$, and the copies at $\frac12$ and at $t(\frac12)=\frac34$ commute.

\emph{Area.} Let $p<q$ in $\mathbb D_m$ and $g=g_{pq}$. The elements $g^{-1}w_p$ and $t^{-1}g^{-1}w_q$ fix $\frac12$ and have length
$O(m)$, so they are represented by words $\omega,\omega'$ of length $O(m)$ in the $h_j$. The words $w_p^{-1}g\omega$ and $w_q^{-1}gt\omega'$ are null
in $F$
and have length $O(m)$, so Guba's quadratic Dehn function~\cite{Guba2006} fills each with $O(m^2)$ cells. Replacing the four
occurrences of $w_p^{\pm1}$ in $[w_pxw_p^{-1},w_qyw_q^{-1}]$ by $(g\omega)^{\pm1}$, and the four of $w_q^{\pm1}$ by $(gt\omega')^{\pm1}$, therefore
costs
$O(m^2)$ cells. Moving $x^{\pm1}$ across $\omega$ and $y^{\pm1}$ across $\omega'$ with the relations $[x,h_j]$ costs $O(m)$ cells and leaves
$g[x,tyt^{-1}]g^{-1}$, a conjugate of a relation. For $p>q$ the commutator is the inverse of the one with the two copies exchanged. So every
area is $O(m^2)$.

\emph{Profile.} Apply Theorem~\ref{thm:criterion} with $N=2^m-1$ and $\Lambda=Cm^2$, where $C\ge1$, and let $l$ be the largest
length of a word of $R$. Given $r$, let $m$ be the largest
integer with $2^mm^2\le r/(10^8C(2l-1))$; then $2^m$ is at least of order $r/\log^2r$, and $\log\sigma_E(r)\ge(2^m-1)/16$. Given
$e$, let $m$ be the largest integer with $2^{m/2}m^2\le1/(CK_\Delta e)$; then $2^m$ is at least of order
$(CK_\Delta)^{-2}e^{-2}/\log^4(1/e)$, and $\log D\ge(2^m-1)\Delta^2/12$.

\emph{No Houghton subgroup.} A torsion element of $\mathcal L$ maps to a torsion element of the torsion-free group $F$, so it lies in
$\bigoplus_{\mathbb D}S_3$, which has exponent six. Every $H_m$ contains a five-cycle. Finally, $\mathcal L$ embeds in $V$ by
Lemma~\ref{lem:gap-coding}.
\end{proof}

Since $\mathcal L$ contains $F$, whose soficity is open, this chunk may have no permutation models at all; if it has any, they are very large.
In particular, either $V$ is not sofic, or its models of this chunk need $2^{c\,r/\log^2r}$ points, close to the largest bound that
Theorem~\ref{thm:criterion} can give. If $F$ is amenable, so is $\mathcal L$, and then this chunk has finite profile. If $F$ is hyperlinear, so
is $\mathcal L$ (Proposition~\ref{prop:thompson-hyperlinear}), and then unitary models of the chunk exist at every defect.

Slofstra's bounds are of the same kind: they bound models that may not exist, since his group is not known to be hyperlinear. Any unitary
models of it need $\log D\gtrsim e^{-c}$ for every $c<\frac12$~\cite[Theorem~1.2 and the remark after it]{Slofstra2018}, and the proof
of~\cite[Corollary~1.3]{Slofstra2018} shows that any permutation models of a suitable chunk need $\log\sigma\gtrsim r^c$ for every $c<\frac14$.
These profiles may be infinite. The group $\mathcal L$ gives $e^{-2}$ and $r$ up to logarithms.

Since $\mathbb D$ is the coset space $F/F_{1/2}$, the group $\mathcal L=\bigl(\bigoplus_{F/F_{1/2}}S_3\bigr)\rtimes F$ has the shape
$\bigl(\bigoplus_{G/H}A\bigr)\rtimes G$ of the generalised wreath products of Kun and Thom~\cite[Theorem~A]{KunThom2026}. They prove such a
product non-sofic for $A=\Z/2$ when $H<G$ is infranormal and not normal, that is, the elements $g$ with $gHg^{-1}\subseteq H$ generate $G$, and
both $H$ and $G$ have property (T). Our input is different: cheap transport of non-abelian copies of $S_3$. Soficity and hyperlinearity pass to
wreath products whose acting group is sofic~\cite{HayesSale2018}. They pass to permutational wreath products when the acting group is sofic
(respectively hyperlinear) and the action is sofic~\cite[Theorems~3.6 and~3.8]{GKEP2025}, \cite{AlekseevBradford2026}. Since $F$ acts on
$\mathbb D$ transitively and faithfully, a sofic action of $F$ on $\mathbb D$ would already make $F$ sofic~\cite[Theorem~2.12]{GKEP2025}, so
these results give nothing here. Proposition~\ref{prop:thompson-hyperlinear} assumes only that $F$ is hyperlinear.

\begin{proposition}[Hyperlinearity transfer]\label{prop:thompson-hyperlinear}
The group $\mathcal L$ is hyperlinear if and only if $F$ is.
\end{proposition}

The proof, in Appendix~\ref{app:hyperlinear}, needs a copy of $S_3$ in the ultrapower $L(F)^\omega$ at $\frac12$ that commutes with the
stabilizer $F_{1/2}$. Near $\frac12$ the stabilizer acts by dilations. The proof therefore places a copy of $S_3$ in each of a sequence of
intervals shrinking to $\frac12$, and selects among them by a random index whose law is invariant under dilation in the limit.
